\documentclass[10pt]{article}

\usepackage[margin=0.72in]{geometry}
\usepackage{amsmath, amssymb}
\usepackage{booktabs}
\usepackage{hyperref}
\usepackage{enumitem}
\setlist{nosep,leftmargin=*}
\setlength{\parskip}{0.22em}
\setlength{\parindent}{0pt}

\title{\textbf{Project Proposal}\\
Submodular Optimization and Greedy Approximation for Distributed Influence Maximization}
\author{Team Members: Yixin Chen, Bichi Zhang, Yaoqin Ye}
\date{CS240: Algorithm Design and Analysis, Spring 2026}

\begin{document}

\maketitle

\section{Topic Selection}

We choose \textbf{Topic 3: Influence Maximization} from the course reference
list. Influence maximization asks how to select a limited number of seed nodes
in a network so that information, behavior, products, or interventions spread
to as many nodes as possible. It is a standard model for viral marketing and
social-network analysis, and it also provides a clean algorithmic connection
between modern networked applications and classical approximation algorithms.

This project focuses on \textbf{submodular optimization and greedy
approximation for distributed influence maximization}. In large networks, no
single machine or agent may know the entire graph. Instead, each local agent
observes part of the graph or estimates local diffusion effects, then
communicates with other agents to approximate a globally effective seed set
under a fixed budget.

\section{Problem Formulation}

Let $G=(V,E)$ be a directed or undirected graph. Each edge $(u,v)$ may have a
propagation probability $p_{uv}$. Given a budget $k$, the goal is to choose a
seed set $S\subseteq V$ with $|S|\leq k$ that maximizes the expected number of
activated nodes:
\[
    \max_{S\subseteq V,\ |S|\leq k} \sigma(S),
\]
where $\sigma(S)$ is the expected influence spread from $S$ under a diffusion
model such as the Independent Cascade model or the Linear Threshold model. The
input is the graph, diffusion parameters, and budget $k$; the output is a seed
set of size at most $k$.

This is a combinatorial optimization problem. Under common diffusion models,
$\sigma(S)$ is monotone and submodular, meaning the marginal gain of adding a
new node decreases as the current seed set grows:
\[
    \sigma(A\cup\{v\})-\sigma(A)
    \geq
    \sigma(B\cup\{v\})-\sigma(B),
    \quad A\subseteq B,\ v\notin B.
\]
Although exact influence maximization is NP-hard, the standard greedy
algorithm achieves a $(1-1/e)$ approximation guarantee for monotone submodular
maximization under a cardinality constraint.

\section{Related Work}

Kempe, Kleinberg, and Tardos introduced the influence maximization problem for
social networks, proved NP-hardness under the Independent Cascade and Linear
Threshold models, and showed how submodular maximization yields a
$(1-1/e)$-approximate greedy algorithm. This will be the main representative
paper for reproduction.

Nemhauser, Wolsey, and Fisher proved the classical approximation guarantee for
maximizing a nonnegative monotone submodular function subject to a cardinality
constraint. Their result provides the theoretical foundation for the greedy
analysis in influence maximization.

Leskovec et al. proposed the Cost-Effective Lazy Forward (CELF) acceleration,
which uses diminishing marginal gains to avoid many repeated influence-spread
computations. We plan to reproduce both standard greedy influence maximization
and CELF-style lazy greedy selection.

\section{Algorithm and Technical Plan}

We plan to implement and compare the following methods:
\begin{itemize}
    \item \textbf{Random baseline:} choose $k$ seed nodes uniformly at random.
    \item \textbf{Degree baseline:} choose the $k$ nodes with highest degree or
    out-degree.
    \item \textbf{Greedy influence maximization:} repeatedly choose the node
    with the largest estimated marginal influence gain.
    \item \textbf{CELF lazy greedy:} maintain upper bounds on candidate
    marginal gains to reduce unnecessary diffusion simulations.
    \item \textbf{Distributed extension:} partition the graph across multiple
    agents, let each agent estimate local marginal gains, and aggregate the
    candidates through limited communication or a consensus-style rule.
\end{itemize}

The main diffusion model will be the Independent Cascade model, with influence
spread estimated by Monte Carlo simulation. Planned datasets include synthetic
Erdos--Renyi graphs, Barabasi--Albert scale-free graphs, Watts--Strogatz
small-world graphs, and one small real network if time permits. We will measure
influence spread, runtime, scalability, and sensitivity to the seed budget $k$
and the number of Monte Carlo trials.

\section{Project Scope and Expected Goals}

The minimum goal is to reproduce the classical behavior of influence
maximization algorithms:
\begin{itemize}
    \item greedy selection should obtain larger influence spread than random
    and degree-based baselines;
    \item CELF should preserve nearly the same solution quality as greedy while
    reducing runtime;
    \item influence spread should increase with $k$, while marginal gains
    decrease because of submodularity;
    \item the distributed approximation should illustrate the tradeoff between
    communication cost and solution quality.
\end{itemize}

Optional extensions include comparing different propagation probabilities,
studying how Monte Carlo sample size affects stability, adding
communication-limited distributed selection, or connecting the results to
distributed optimization with local information and global objectives.

\section{Team Information}

\begin{tabular}{@{}lll@{}}
\toprule
Member & Responsibility & Expected Contribution \\
\midrule
Yixin Chen & Modeling and theory & Problem formulation and related work \\
Bichi Zhang & Algorithms & Greedy, CELF, and complexity analysis \\
Yaoqin Ye & Experiments & Implementation, plots, and distributed extension \\
\bottomrule
\end{tabular}

\end{document}
